\section{Architecture}

\subsection{Layers}

The testbed implements five layers, each emitting events into one normalized
schema:

\begin{enumerate}
\item \textbf{Spacecraft.} Circular-orbit propagation, power/thermal/attitude
dynamics with physically bounded rates, a thirteen-command catalogue with a
\emph{critical} flag, SDLS tag verification and an anti-replay window.
\item \textbf{Ground station.} Pass geometry, an elevation-dependent frame error
model, integrity verification of downlinked frames, and normalization of every
observation into the event schema.
\item \textbf{Mission-control API.} Authentication with optional second factor,
role-based command allowlists, per-principal rate limiting, and a pass-window
feasibility check. Frames are signed here before transmission.
\item \textbf{Cloud.} A stream carrying normalized events, an encrypted object
store for telemetry products, an audit trail covering both management and object
data events, a streaming detection function, findings aggregation, metrics and
alarms.
\item \textbf{Response.} A containment function subscribed to high and critical
findings, implementing the containment step --- and only the containment step ---
of the incident-response playbooks.
\end{enumerate}

\subsection{Two paths, two latencies}

The architecture has exactly two detection paths and their latencies differ by an
order of magnitude:

\begin{description}
\item[Stream path.] Ground-segment observations (authentication, telecommands,
frame integrity, telemetry samples) reach the detection function through the
stream within seconds. Its bound is the consumer's batching window.
\item[Control-plane path.] Cloud API calls reach the same function through the
audit trail and the event bus. Its bound is trail delivery, which is minutes, not
seconds.
\end{description}

We model these as a per-rule \emph{tier} (\texttt{stream} or \texttt{batch}) with
a configurable latency added to every alert timestamp. Reporting a detection
latency that excludes pipeline delay would flatter the architecture; making the
delay an explicit, swept parameter (\S\ref{sec:results}) makes the claim honest
and lets a reader substitute latencies measured in their own account.

\subsection{Normalized event schema}

Every layer emits the same record: timestamp, source, event type, actor, target,
outcome, source address, and a rule-specific attribute map. The schema is a small
subset of the common event models used in practice --- enough to write portable
detections, small enough that its data dictionary fits on a page.

Each event additionally carries a ground-truth label (\texttt{benign} or an attack
identifier). The label is written by the scenario, never read by a rule --- a
property enforced by a test that fails if the string appears anywhere in the rule
module --- and used only for scoring.

\subsection{Deployment}

The cloud instantiation is Terraform for a stream, an encrypted archive with
lifecycle transitions, an audit trail with object-level data events on the
archive, two functions, an event bus, findings aggregation, and least-privilege
roles in which the detector can observe but not act and the responder can act
only through a narrow set of reversible containment operations. The detection
function imports the same rule package the experiments run, shipped as a layer:
a single implementation is what allows results measured offline to be claimed for
the deployed system.

\subsection{Automated response, deliberately constrained}

Only containment is automated, and only with reversible actions. A ground segment
controls a vehicle that cannot be taken offline for maintenance; an automated
action that safes a spacecraft in response to a false positive causes a real
mission outage, and the response becomes the incident. The most aggressive action
the automation can take is to stop the \emph{ground segment} from transmitting.
Nothing in the response path commands the spacecraft.
